%% kapa: exam=1 round=20 year=2009 period=2 form=A subject=appraisal_law subject_ko=부동산관계법규 seq=15 number=55
%% source: Q-Net 감정평가사 기출문제 (cst00309) artlSeq=5024178 file=extracted/1008/5024178/2009년 제20회 감정평가사 1차시험 기출문제/2교시.hwp
%% conversion: kordoc 4.24.51 -> tools/gen1.py (auto); review=proofread-sonnet; flags=none
\begin{question}{55}
국토의 계획 및 이용에 관한 법령상 도시계획시설사업 시행자의 수용권 또는 사용권의 범위에 관한 설명으로 옳지 \uline{않은} 것은?
\begin{choices}
  \item 도시계획시설사업 시행자는 도시계획시설사업에 필요한 토지의 지상권을 수용할 수 있다.
  \item 도시계획시설사업 시행자는 도시계획시설사업에 필요한 건축물의 전세권을 수용할 수 있다.
  \item 도시계획시설사업 시행자는 도시계획시설사업에 필요한 토지에 정착된 나무를 수용할 수 있다.
  \item 도시계획시설사업 시행자는 도시계획시설사업 시행을 위하여 특히 필요하다고 인정되는 경우 도시계획시설에 인접한 건축물을 일시 사용할 수 있다.
  \item 도시계획시설사업 시행자는 도시계획시설사업 시행을 위하여 특히 필요하다고 인정되는 경우 도시계획시설에 인접한 토지를 수용할 수 있다.
\end{choices}
\end{question}
\answer{5}
